\documentclass[10pt,a4paper]{amsart}
\usepackage[margin=19mm]{geometry}
\usepackage{amsmath,amssymb,mathtools}
\usepackage[scr=rsfs,cal=euler]{mathalfa}
\usepackage{xfrac}
\usepackage[unicode,hidelinks,bookmarks=false]{hyperref}
\newcommand{\ie}{i.e.\ }
\newcommand{\cf}{cf.\ }
\newcommand{\resp}{respectively\ }
\newcommand{\Subsection}{\subsection}
\providecommand{\nobreakdash}{\nobreak\hbox{-}}
\setlength{\emergencystretch}{2em}
\multlinegap0pt
\usepackage{float}
\usepackage{cancel}
\usepackage{tikz}
\usepackage{mathtools}
\usepackage{bm}
\usepackage[shortlabels]{enumitem}
\newtheorem{theorem}{Theorem}
\newtheorem{lemma}{Lemma}
\newtheorem{prop}{Proposition}
\newcommand{\C}{\mathbb{C}}
\DeclareMathOperator{\Imz}{Im}
\newcommand{\N}{\mathbb{N}}
\newcommand{\R}{\mathbb{R}}
\DeclareMathOperator{\arcsinh}{arcsinh}
\title{Complex Methods in the Asymptotics of M\"obius energy integrals of helix curves}
\author{Max Lipton}
\date{Source: arXiv:2605.12815v2 (CC BY 4.0)}
\begin{document}
\maketitle

\noindent\emph{Source and licence.} Complex Methods in the Asymptotics of M\"obius energy integrals of helix curves. Version arXiv:2605.12815v2 (CC BY 4.0), distributed by arXiv under the Creative Commons Attribution 4.0 International licence, http://creativecommons.org/licenses/by/4.0/, as recorded in the arXiv licence metadata for this exact version and shown on the abstract page https://arxiv.org/abs/2605.12815v2. Full licence notice: \texttt{licenses/arxiv-2605.12815-CC-BY-4.0.txt}.\par\smallskip

\noindent\textit{Editorial scope.} Counted unit: the article's prop jtransferthm (source0.tex lines 1032--1038). The article's complete original proof of it is delivered with it (source0.tex lines 1040--1117, one \texttt{proof} environment of 78 lines). No mathematics is written, repaired or re-derived: the statement and the proof are the article's own lines, in the article's own notation, and no retained line is substituted or re-flowed.  Retained uncounted as the source-local context the proof needs: lines 283--295; lines 549--551; lines 791--793; lines 796--801; lines 798--800; lines 997--1003.  Nothing was omitted from any retained span, so the package carries no omission record.  No retained line contains a citation, so the wrapper carries no bibliography and no bibliography entry.  No retained line contains a figure environment, an image-inclusion command or drawing code, so no image is delivered and the package declares no asset dependency.  Editorial formatting only, in addition: an AMS \texttt{amsart} preamble that restates the article's own declarations and macros used by the retained lines, namely the environments \texttt{theorem}, \texttt{lemma}, \texttt{prop} on separate counters, together with the standard packages those lines need.  The retained environments print in this document as Theorem 1, Lemma 1, Proposition 1 in order of appearance, whereas the article prints the same items with the numbering of its own full document.  The complete native source of the article as distributed in the arXiv e-print for this exact version is bundled unchanged as \texttt{sources/200455/source0.tex}.
\begin{theorem}[Transfer Theorem] \label{transferthm}
    Let $F,G: (0,\infty) \to \C$ be complex functions which admit absolutely convergent series representations. That is, we have one parameter families of sequences of complex numbers $a_k(\rho)$ and $b_k(\rho)$ such that for all $\rho > 0$,
     \begin{align*}
    F(\rho) &\coloneq \sum_{k=1}^{\infty} a_k(\rho), \quad
    G(\rho) \coloneq \sum_{k=1}^{\infty} b_k(\rho),
\end{align*}
where all the series are absolutely convergent. Suppose there exists $c>0$ and $\rho_0 > 0$ such that for all $0 < \rho < \rho_0$ and $k \in \N$, the following conditions hold:
\begin{enumerate}
    \item (Sign condition) The imaginary parts of all $b_k(\rho)$ share the same sign and are nonzero.
    \item (Non-tangency) $\left| \Imz b_k(\rho) \right| \geq c|b_k(\rho)|.$
\end{enumerate}
Then if $F(\rho) \sim G(\rho)$ as $\rho \to 0$, we have that $\Imz F(\rho) \sim \Imz G(\rho)$ as $\rho \to 0$.
\end{theorem}

\begin{lemma}\label{arcsinhu_overu}
For all $u \in \R / \{0\}$, we have $0 < \frac{\arcsinh u}{u} < 1$.
\end{lemma}

\begin{align}
\left| E'_\rho(w_k(\rho)) \right| &\geq 4k\pi\rho \sqrt{(k\pi\rho)^2 + 1}.    \label{derho_wk_lowerbound}
\end{align}

\begin{prop}\label{zkrk_bound_thm}
    For all $\rho > 0$ sufficiently small and for all $k \in \N$, we have
    \begin{align}
        |z_k(\rho) - w_k(\rho)| &\leq 2\sqrt{5} \frac{\rho \arcsinh k\pi\rho}{\sqrt{(k\pi\rho)^2 + 1}}. \label{zkrk_bound}
    \end{align}
\end{prop}

    \begin{align}
        |z_k(\rho) - w_k(\rho)| &\leq 2\sqrt{5} \frac{\rho \arcsinh k\pi\rho}{\sqrt{(k\pi\rho)^2 + 1}}. 
    \end{align}

\begin{prop}\label{jtildetransferthm}
The function $\tilde{J}$ satisfies the condition of the transfer theorem, Theorem \ref{transferthm}. In fact, for all $\rho > 0$ and $k \in \N$,
\begin{enumerate}
    \item $\Imz \frac{1}{E'_\rho(w_k(\rho))} < 0$, and
    \item $\left| \Imz \frac{1}{E'_\rho(w_k(\rho))} \right| \geq \frac{\pi}{\sqrt{\pi^2 + 1}} \left| \frac{1}{E'_\rho(w_k(\rho))} \right|$.
\end{enumerate}  
\end{prop}

\begin{prop}\label{jtransferthm}
The function $J$ satisfies the condition of the transfer theorem, Theorem \ref{transferthm}. There exists $c > 0$ and $\rho_0 > 0$ such that for all $0 < \rho < \rho_0$ and $k \in \N$,
\begin{enumerate}
    \item $\Imz \frac{1}{E'_\rho(z_k(\rho))} < 0$, and
    \item $\left| \Imz \frac{1}{E'_\rho(z_k(\rho))} \right| \geq c\left| \frac{1}{E'_\rho(z_k(\rho))} \right|$.
\end{enumerate}  
\end{prop}

\begin{proof}
 We fix $k \in \N$, and begin with the assumption that $\rho > 0$ is sufficiently small so that the estimates on $|z_k - w_k|$ from Prop. \ref{zkrk_bound_thm} apply. Let $C_R = 2\sqrt{5}$ be the bounding constant from that proposition. The $R$ stands for Rouch\'e.

We begin with an estimate of $|E'_\rho(z_k) - E'_\rho(w_k)|$. By the fundamental theorem of calculus,
\begin{align*}
    E'_\rho(z_k) - E'_\rho(w_k) = (z_k-w_k)\int_0^1 E''_\rho(w_k + t(z_k-w_k))dt. 
\end{align*}
Let $\zeta = \zeta_k = \zeta_k(\rho) \coloneq z_k - w_k$, so we have that 
\begin{align*}
    |\zeta| &\leq C_R \frac{\rho \arcsinh k\pi\rho}{\sqrt{(k\pi\rho)^2 + 1}} \\
    &\leq C_R \rho,
\end{align*}
where we have used the fact $ \arcsinh u < \sqrt{u^2 + 1} $ for all $u>0$. Let us further assume $\rho < \frac{1}{C_R}$, which ensures $|\zeta| < 1$. This allows us to apply the same argument used to bound $M = \sup\limits_{|\zeta - w_k| \leq r_k} |E''_\rho(\zeta)|$ from Prop. \ref{zkrk_bound_thm} to see that for all $t \in [0,1]$,
\begin{align*}
    \left| E''_\rho(w_k + t\zeta)\right| &\leq 11((k\pi\rho)^2 + 1).
\end{align*}
Therefore,
\begin{subequations}
\begin{align}
    |E'_\rho(z_k) - E'_\rho(w_k)| &\leq 11|\zeta|((k\pi\rho)^2 + 1)  \nonumber\\
    &\leq 11C_R \rho \arcsinh( k\pi\rho) \sqrt{(k\pi\rho)^2 + 1}, \label{dEzkwk}
\end{align}
where we applied \eqref{zkrk_bound}.
Now apply \eqref{derho_wk_lowerbound} and Lemma \ref{arcsinhu_overu} to yield
\begin{align}
  \left| \frac{E'_\rho(z_k)- E'_\rho(w_k)}{E'_\rho(w_k)} \right| &=    \frac{|E'_\rho(z_k) - E'_\rho(w_k)|}{|E'_\rho(w_k)|}\nonumber \\
    &\leq \frac{11C_R}{4} \rho. \label{dEzkwk_asymp}
\end{align}
\end{subequations}
The estimate \eqref{dEzkwk_asymp} allows us to write
\begin{subequations}
\begin{align}
    E'_\rho(z_k) &= E'_\rho(w_k)(1 + \eta), \label{dErho1+eta}
\end{align}
for some $\eta \in \C$ such that $|\eta| \leq \frac{11C_R}{4}\rho$. In turn, we have
\begin{align}
    \frac{1}{E'_\rho(z_k)} &= \frac{1}{E_\rho'(w_k)(1+\eta)} \nonumber \\
    &= \frac{1}{E_\rho'(w_k)} \left(1 - \frac{\eta}{1+\eta} \right). \label{dErhoinverse1+eta}
\end{align}
Now let us assume $\rho$ is small enough to guarantee $|\eta| \leq \frac{1}{2}$. Then the reverse triangle inequality yields $|1 + \eta| \geq | |1| - |\eta|| = 1 - |\eta|$. Therefore we have the controlling inequality,
\begin{align}
    \left| \frac{\eta}{1 + \eta} \right| &\leq \frac{|\eta|}{1 - |\eta|} \nonumber\\
    &\leq \frac{ \frac{11}{4}C_R \rho}{1 - |\eta|} \nonumber\\
    &\leq \frac{11}{2}C_R \rho. \label{etacontrol}
\end{align}
\end{subequations}
Now see that by the sign and non-tangency conditions shown in Prop. \ref{jtildetransferthm} and the trivial estimate $|\Imz z| \leq |z|$,
\begin{align*}
    \Imz \frac{1}{E'_\rho(z_k)} &= \Imz \left[ \frac{1}{E'_\rho(w_k)} \left(1 + \frac{-\eta}{1+\eta} \right)\right] \\
    &= \Imz \frac{1}{E'_\rho(w_k)} + \Imz \left[ \frac{-\eta}{E'_\rho(w_k)(1+\eta)}\right] \\
    &\leq -\frac{\pi}{\sqrt{\pi^2 + 1}}\left| \frac{1}{E'_\rho(w_k)} \right| +  \left| \frac{\eta}{1+\eta} \right|\left| \frac{1}{E'_\rho(w_k)} \right|. \\
\end{align*}
Thus, we can guarantee $\Imz \frac{1}{E'_\rho(z_k)} < 0$ by assuming $\rho$ is small enough to guarantee $\left| \frac{\eta}{1+\eta} \right| < \frac{\pi}{\sqrt{\pi^2 + 1}}$.
To complete this proof, it now suffices to find $c_1,c_2 > 0$ such that
\begin{align*}
    \left| \Imz \frac{1}{E'_\rho(z_k)} \right| &\geq c_1 \left| \Imz \frac{1}{E'_\rho(w_k)} \right|
\end{align*}
and
\begin{align*}
    \left| \frac{1}{E'_\rho(w_k)} \right| \geq c_2 \left| \frac{1}{E'_\rho(z_k)} \right|,
\end{align*}
joining this chain by the non-tangency of $\frac{1}{E'_\rho(w_k)}$.
We have just seen
\begin{align*}
    \left| \Imz \frac{1}{E'_\rho(z_k)} \right| &= -\Imz \frac{1}{E'_\rho(z_k)} \\
    &\geq \left( \frac{\pi}{\sqrt{\pi^2 + 1}} - \left| \frac{\eta}{1+\eta} \right| \right) \left| \frac{1}{E'_\rho(w_k)} \right| \\
    &\geq \left( \frac{\pi}{\sqrt{\pi^2 + 1}} - \left| \frac{\eta}{1+\eta} \right| \right) \left| \Imz \frac{1}{E'_\rho(w_k)} \right|.
\end{align*}
Now we shrink $\rho$ further to ensure $|\frac{\eta}{1+\eta}| < \frac{1}{2}$, which allows us to take $c_1 = \frac{\pi}{\sqrt{\pi^2 + 1}} - \frac{1}{2} > 0$.
 
Furthermore, from \eqref{dErhoinverse1+eta},
\begin{align*}
\left| \frac{1}{E'_\rho(z_k)} \right| &= \left| 1 - \frac{\eta}{1+\eta} \right|\left| \frac{1}{E'_\rho(w_k)} \right|\\
&\leq \left(1 + \left|\frac{\eta}{1+\eta}\right| \right) \left| \frac{1}{E'_\rho(w_k)} \right| \\
&\leq \frac{3}{2}\left|\frac{1}{E'_\rho(w_k)} \right|,
\end{align*}
allowing us to take $c_2 = \frac{2}{3}$.
\end{proof}

\end{document}
